
\documentclass[a4paper,landscape]{article}
\usepackage{longtable}
\usepackage[utf8]{inputenc}

\begin{document}
\begin{longtable}{|p{5cm}|p{5cm}|p{5cm}|c|c|c|c|}
\hline
\textbf{Original (RU)} & \textbf{Translation (EN)} & \textbf{Summary of Evaluation} & \textbf{Fluency} & \textbf{Faithfulness} & \textbf{Tone/Style} & \textbf{Combined Total} \\
\hline
\endfirsthead
\hline
\endfoot
только знакомя нас с героями Достоевский уже подготавливает почву для это их жуткой ночной встречи Кириллов у него молодой человек лет около 27 верховенский почти зеркально молодой человек лет 27 или около это нарочитое сходство Разумеется только подчеркивает Сколь глубоки между ними & is already preparing the ground for their terrible night meeting, Kirillov, he has a young man of about 27 years old, almost Verkhovensky a mirror-image young man of about 27 years old or so, this deliberate resemblance, of course, only emphasizes how deep the differences are between them. Kirillov is a bright representative of & Fluency awkward but understandable; Faithfulness high; Tone slightly lost. & 0.5 & 1.0 & 0.5 & 2.0 \\
\hline
XIX века часто пыталась решить как человеку жить После отказа от Бога допустим Бога больше нет Допустим мы в него больше не верим мы его как выразился бы верховенски раскосировали Но ведь потребность чем-то заполнить пустоту на его месте она-то осталась и & a person can live. After abandoning God, let’s say there is no God anymore Suppose we no longer believe in him, we have, as Verkhovenski put it, deformed him , but the need to fill the void with something in his place remains, and & Fluency is clear; Faithfulness is good but loses some nuance; Tone and style mostly intact. & 1.0 & 0.5 & 1.0 & 2.5 \\
\hline
если Бога Нет Значит я могу занять его место помните Гребенщиков пел ты чувствуешь сквозняк от того что это место свободно Кириллов чувствует тот самый сквозняк и он формулирует страшную безумную мысль в & place , remember Grebenshchikov sang, you feel a draft from the fact that this place is free Kirillov feels that same draft and he formulates a terrible crazy thought in which Then he himself falls in love, if & Fluency suffers due to broken structure; Faithfulness somewhat lost due to misinterpretation; Tone is entirely lost. & 0.5 & 0.5 & 0.0 & 1.0 \\
\hline
человека Богом Так ведь самоубийц уже миллионы были говорят ему но все не Затем все со страхом и не для того возражает Кириллов и этот его Безумный подход удивительным образом напоминает один эпизод из Дон & a human God. So after all, there have been millions of suicides already, they tell him, but still nothing. Then Kirillov objects with fear and for the wrong reasons, and this Insane approach of his surprisingly resembles one episode from Don & Fluency is good, some slight awkwardness; Faithfulness is mostly preserved; Style and tone could be better. & 1.0 & 0.5 & 0.5 & 2.0 \\
\hline
и этот его Безумный подход удивительным образом напоминает один эпизод из Дон Кихота вспомните когда дерзновенные и далеко говорит Санчо панси что планирует подражать славным рыцарям прошлого и Сойти с ума & and this Insane approach of his surprisingly resembles one episode from Don Quixote, remember when the daring and far away Sancho says to Pansi that he plans to imitate the glorious knights of the past and to go crazy, & Fluency is clear; Faithfulness is impacted by misattribution of speaker; Tone is preserved. & 1.0 & 0.0 & 0.5 & 1.5 \\
\hline
паузы сразу ба Да я что-то примечаю там на окне на тарелке Он подходит к окну Ого вареная с рисом курица но почему же до сих пор непочатая стало быть мы находимся в таком настроении Духа что даже и курицу я ел перебивает его Кириллов и не ваше дело молчите о конечно И притом все равно но для меня то оно теперь не равно вообразите совсем & a pause immediately ba Yes, I notice something there on the window on a plate He approaches the window Wow, chicken boiled with rice but why is it still unfinished, so we are in such a mood of the Spirit that even I ate the chicken, Kirillov interrupts him and it’s none of your business keep quiet about of course And besides, it’s all the same, but for me it’s not the same now imagine I almost didn’t have lunch at all and therefore if now this & Fluency is good; Faithfulness excellent; Tone/style preserved well. & 1.0 & 1.0 & 1.0 & 3.0 \\
\hline
и Обратите внимание как формулирует свой ответ Кириллов Ешьте если можете верховенский может еще как может за себя из-за того парня может вот благодарю говорит он а потом и чаю и дальше следует что-то еще более жуткое косточка & Kirillov formulates his answer Eat if you can Verkhovensky can do as well as he can for himself because of that guy, maybe I thank him, he says and then tea and then something even more creepy follows the bone on the cake of this scene Verkhovensky instantly & Fluency is fine; Faithfulness is decent but lacks depth in metaphor; Tone/style works well. & 1.0 & 0.5 & 1.0 & 2.5 \\
\hline
из-за того парня может вот благодарю говорит он а потом и чаю и дальше следует что-то еще более жуткое косточка на торте этой сцены верховенский мигом устроился за столом на другом конце дивана из чрезвычайную жадностью & because of that guy, maybe I thank him, he says and then tea and then something even more creepy follows the bone on the cake of this scene Verkhovensky instantly settled down at the table at the other end of the sofa from he pounced on the food with extreme greed, & Fluency is fine; Faithfulness is good but a slight loss of nuance; Tone/style mostly retained. & 1.0 & 1.0 & 0.5 & 2.5 \\
\hline
он ест жадно с аппетитом с азартом ведь сидящий в метре от него человек сейчас застрелится но и курица стала быть вакантно заметьте в одном этом абзаце верховенский забирает у Кириллова & the man sitting a meter away from him is about to shoot himself, but the chicken also began to be vacant, notice in this one paragraph, Verkhovensky takes away from Kirillov absolutely all his freedom, his last & Fluency is smooth; Faithfulness good but a slight loss of nuance; Tone/style mostly retained. & 1.0 & 0.5 & 1.0 & 2.5 \\
\hline
Людоед люди для него расходный материал питающий его жалкие амбиции и жуткие планы Стоит ли счастье человечество слезинки ребенка спрашивает Иван карамазов А ребеночек с гарниром или так без всего уточняет Петр и Кириллов этот странный Кириллов Да со & people are consumable for him, feeding his pathetic ambitions and terrible plans Is happiness worth humanity’s tears of a child asks Ivan Karamazov And a baby with a side dish or so without anything, Peter and Kirillov clarify this strange Kirillov Yes, with & Fluency is clear; Faithfulness is solid; Tone/style is appropriate but lacks some nuance. & 1.0 & 0.5 & 1.0 & 2.5 \\
\hline
он поедает но и его Кириллова он сейчас съест верховенский Барабек скушал 40 человек и корову и Быка и Кривого мясника и телегу и дугу и метлу и кочергу скушал курицу с гарниром и покончил с этим миром и весь этот клубок противоречий страхов и судеб Достоевский показывает через & he eat chicken, but he will also eat his Kirillov , Verkhovensky Barabek ate 40 people and the cow and the Bull and the Crooked Butcher and the cart and the arc and the broom and the poker ate the chicken with a side dish and ended with this world and this whole tangle of contradictions of fears and destinies Dostoevsky shows through & Fluency is smooth; Faithfulness lost due to misheard content; Tone/style is entirely absent. & 1.0 & 0.0 & 0.0 & 1.0 \\
\hline
Я хочу это что Марина пойдёт & I want this, that Marina will go, & Fluency is fragmented and awkward; Meaning is lost due to misheard content; Tone/style is missing. & 0.5 & 0.0 & 0.0 & 0.5 \\
\hline
остается только оговориться что хотят духу Дон Кихота сервантеса этот эпизод соответствует вполне но букве не полностью в тексте романа именно вот такой сцены с курицей не было Однако это в конце концов легко объяснимо Сервантес & all that remains is to make a reservation that they want the spirit of Don Cervantes's Quixote this episode corresponds completely but not completely to the letter in the text of the novel there was no such scene with the chicken. However, this is, in the end, easily explainable. Cervantes & Fluency is strong; Faithfulness is solid; Tone/style captured but speaker attribution is wrong. & 1.0 & 1.0 & 0.5 & 2.5 \\
\hline
Total & & & 11.5 & 7.0 & 7.5 & 26.0 \\
\end{longtable}
\end{document}
